\chapter{Asymptotic Cones and Sharafutdinov Retraction}
\label{app:chow-iii-asymptotic-cones-sharafutdinov}

\providecommand{\M}{\mathcal M}
\providecommand{\N}{\mathcal N}
\providecommand{\R}{\mathbb R}
\providecommand{\Ray}{\operatorname{Ray}}
\providecommand{\sect}{\operatorname{sect}}
\providecommand{\inj}{\operatorname{inj}}
\providecommand{\dhat}{\widehat d}
\providecommand{\ddinf}{\widetilde d_\infty}
\providecommand{\Shar}{\operatorname{Shar}}

\section{Sharafutdinov retraction}

\begin{definition}[Rays and Busemann functions]
\label{def:chow-iii-i1-rays-busemann}
\source{chow-et-al-ricci-flow-part-iii:page-0479}
\dcref{appi:1.1}\group{Busemann geometry}
For a noncompact Riemannian manifold $(\M^n,g)$, let $\Ray_{\M}$ be the set of
unit-speed geodesic rays, and let $\Ray_{\M}(p)$ denote those starting at $p$.
For $\gamma\in\Ray_{\M}$ define
\[
b_\gamma(x)=\lim_{s\to\infty}(s-d(\gamma(s),x)),\qquad
b_p=\sup_{\gamma\in\Ray_{\M}(p)}b_\gamma.
\]
\end{definition}

\begin{lemma}[Busemann bounds]
\label{lem:chow-iii-i2-i3-busemann-bounds}
\source{chow-et-al-ricci-flow-part-iii:page-0479}
\uses{def:chow-iii-i1-rays-busemann}
\dcref{appi:1.2,appi:1.3}\group{Busemann geometry}

If $\gamma$ starts at $p$, then $|b_\gamma(x)|\le d(x,p)$ and
$|b_\gamma(x)-b_\gamma(y)|\le d(x,y)$.  The same two inequalities hold for
$b_p$.
\end{lemma}
\begin{proof}
\sketch
The triangle inequality bounds each finite expression
$s-d(\gamma(s),x)$ by $\pm d(x,p)$ and bounds its difference at $x,y$ by
$d(x,y)$.  Pass to the limit and then take the supremum over rays from $p$.
\end{proof}

\begin{lemma}[Long geodesics are asymptotic to rays]
\label{lem:chow-iii-i4-ray-angle}
\source{chow-et-al-ricci-flow-part-iii:page-0480}
\uses{def:chow-iii-i1-rays-busemann}
\dcref{appi:1.4}\group{Busemann geometry}

If $\theta(r)$ is the supremal angle at $p$ between an arbitrary minimizing
segment of length at least $r$ and the closest ray from $p$, then
$\lim_{r\to\infty}\theta(r)=0$.
\end{lemma}

\begin{definition}[Geodesic half-spaces and Busemann sublevels]
\label{def:chow-iii-i5-halfspaces}
\source{chow-et-al-ricci-flow-part-iii:page-0480}
\uses{def:chow-iii-i1-rays-busemann}
\dcref{appi:1.5}\group{Busemann geometry}

For $\gamma\in\Ray_{\M}(p)$ set
\[
\mathbb B_\gamma=\bigcup_{s>0}B(\gamma(s),s),\qquad
\mathbb H_\gamma=\M\setminus\mathbb B_\gamma,
\]
and write $\gamma_s(u)=\gamma(u+s)$.  Define
\[
C_s(p)=\bigcap_{\gamma\in\Ray_{\M}(p)}\mathbb H_{\gamma_s}.
\]
\end{definition}

\begin{lemma}[Half-space and sublevel identities]
\label{lem:chow-iii-i6-i9-halfspace-sublevels}
\source{chow-et-al-ricci-flow-part-iii:page-0480,chow-et-al-ricci-flow-part-iii:page-0481}
\uses{def:chow-iii-i5-halfspaces, lem:chow-iii-i2-i3-busemann-bounds}
\dcref{appi:1.6,appi:1.7,appi:1.8,appi:1.9}\group{Busemann geometry}

For $0<s<t$, $\mathbb B_{\gamma_s}$ is the $(t-s)$-neighborhood of
$\mathbb B_{\gamma_t}$ and $\mathbb H_\gamma=\{b_\gamma\le0\}$.  Moreover,
\[
C_s(p)=\{x:b_p(x)\le s\},\quad \partial C_s(p)=\{b_p=s\},\quad
C_s(p)\subset C_t(p)\ (s\le t),\quad \overline B(p,s)\subset C_s(p).
\]
\end{lemma}
\begin{proof}
\sketch
The first identity follows by applying the triangle inequality along the ray;
the remaining identities follow by taking intersections over rays and using
the definition of $b_p$.
\end{proof}

\begin{example}[Convexity examples]
\label{ex:chow-iii-i10-i11-convexity}
\source{chow-et-al-ricci-flow-part-iii:page-0481}
\dcref{appi:1.10,appi:1.11}\group{Busemann geometry}
On the round sphere, a nonempty totally convex subset is the whole sphere,
while a strict spherical cap is strongly convex and the open hemisphere is not.
A convex subset of Euclidean space is totally convex.
\end{example}

\begin{proposition}[Busemann convexity and half-space convexity]
\label{prop:chow-iii-i12-i17-busemann-convexity}
\source{chow-et-al-ricci-flow-part-iii:page-0482,chow-et-al-ricci-flow-part-iii:page-0483}
\uses{lem:chow-iii-i4-ray-angle, lem:chow-iii-i6-i9-halfspace-sublevels}
\dcref{appi:1.12,appi:1.13,appi:1.15,appi:1.16,appi:1.17}\group{Nonnegative curvature and Busemann functions}

If $\sect(g)\ge0$, every $b_\gamma$ and $b_p$ is convex, and if
$\operatorname{Rc}\ge0$ then $b_\gamma$ is distributionally subharmonic.  In
addition,
\[
b_p(x)\ge d(x,p)(1-\theta(d(x,p))),
\]
closed left half-spaces are totally convex, and every $C_s(p)$ is compact and
totally convex.  For $s<t$,
\[
C_s(p)=\{x\in C_t(p):d(x,\partial C_t(p))\ge t-s\}.
\]
\end{proposition}

\begin{theorem}[Cheeger--Gromoll splitting]
\label{thm:chow-iii-i18-splitting}
\source{chow-et-al-ricci-flow-part-iii:page-0483}
\uses{prop:chow-iii-i12-i17-busemann-convexity}
\dcref{appi:1.18}\group{Nonnegative curvature and Busemann functions}

If $(\M^n,g)$ is complete with $\operatorname{Rc}\ge0$ and contains a geodesic
line, then $(\M,g)$ is isometric to $\R\times(\N^{n-1},h)$ with
$\operatorname{Rc}_h\ge0$.
\end{theorem}

\begin{definition}[Soul]
\label{def:chow-iii-i20-i21-soul}
\source{chow-et-al-ricci-flow-part-iii:page-0483,chow-et-al-ricci-flow-part-iii:page-0484}
\dcref{appi:1.20,appi:1.21}\group{Soul theory}
A submanifold $S\subset\N$ is totally geodesic when its second fundamental
form vanishes; its normal bundle is
$\nu(S)=\{V\in T_p\N:p\in S,\ \langle V,T_pS\rangle=0\}$.  A soul is a
compact, totally convex, totally geodesic submanifold.
\end{definition}

\begin{theorem}[Soul theorem and soul conjecture]
\label{thm:chow-iii-i22-i23-soul}
\source{chow-et-al-ricci-flow-part-iii:page-0484}
\uses{def:chow-iii-i20-i21-soul, prop:chow-iii-i12-i17-busemann-convexity}
\dcref{appi:1.22,appi:1.23}\group{Soul theory}

If $\sect(g)\ge0$, a complete noncompact $\M$ has a soul $S$ and is
diffeomorphic to $\nu(S)$; any two souls are isometric.  If $\sect(g)>0$,
every soul is a point.  More generally, nonnegative sectional curvature with
positive sectional curvature somewhere forces the soul to be a point.
\end{theorem}

\begin{theorem}[Sharafutdinov retraction]
\label{thm:chow-iii-i24-i25-sharafutdinov}
\source{chow-et-al-ricci-flow-part-iii:page-0484,chow-et-al-ricci-flow-part-iii:page-0485,chow-et-al-ricci-flow-part-iii:page-0486}
\uses{prop:chow-iii-i12-i17-busemann-convexity, def:chow-iii-i20-i21-soul}
\dcref{appi:1.24,appi:1.25}\group{Soul theory}

For $a_2<a_1\le(-b_p)_{\max}$ there is a distance-nonincreasing map
\[
\Phi:(-b_p)^{-1}(a_2)\longrightarrow(-b_p)^{-1}(a_1).
\]
These level-set maps assemble into a continuous distance-nonincreasing
retraction $\Shar:\M\to C_{\max}(p)$, equal to the identity on $C_{\max}(p)$,
where $C_{\max}(p)=\{x:-b_p(x)=(-b_p)_{\max}\}$.
\end{theorem}

\section{Existence of asymptotic cones}

\begin{definition}[Ideal-boundary angle]
\label{def:chow-iii-i27-ideal-angle}
\source{chow-et-al-ricci-flow-part-iii:page-0486}
\uses{def:chow-iii-i1-rays-busemann}
\dcref{appi:eq-1.7}\group{Asymptotic cones}

For rays $\gamma_1,\gamma_2\in\Ray_{\M}(p)$ define
\[
\ddinf(\gamma_1,\gamma_2)=\lim_{s,t\to\infty}
\zeta_{\gamma_1(s)p\gamma_2(t)}\in[0,\pi],
\]
where $\zeta$ is the Euclidean comparison angle.
\end{definition}

\begin{lemma}[Pseudo-metric on rays]
\label{lem:chow-iii-i27-pseudometric}
\source{chow-et-al-ricci-flow-part-iii:page-0486,chow-et-al-ricci-flow-part-iii:page-0487}
\uses{def:chow-iii-i27-ideal-angle}
\dcref{appi:1.27}\group{Asymptotic cones}

For $a,b>0$,
\[
\lim_{t\to\infty}\frac{d(\gamma_1(at),\gamma_2(bt))}{t}
=\bigl(a^2+b^2-2ab\cos\ddinf(\gamma_1,\gamma_2)\bigr)^{1/2}.
\]
The function $\ddinf$ is a pseudometric, and the ideal boundary is
$\M(\infty)=\Ray_{\M}(p)/\{\ddinf=0\}$.
\end{lemma}

\begin{theorem}[Asymptotic cone]
\label{thm:chow-iii-i26-asymptotic-cone}
\source{chow-et-al-ricci-flow-part-iii:page-0486,chow-et-al-ricci-flow-part-iii:page-0488,chow-et-al-ricci-flow-part-iii:page-0489}
\uses{lem:chow-iii-i27-pseudometric}
\dcref{appi:1.26}\group{Asymptotic cones}

If $(\M^n,g)$ is complete, noncompact, and has $\sect(g)\ge0$, then the
pointed rescalings $(\M,\lambda^2g,p)$ converge as $\lambda\downarrow0$ to the
Euclidean metric cone $\operatorname{Cone}(\M(\infty),\ddinf)$.
\end{theorem}

\section{Distance spheres and critical points}

\begin{definition}[Radial intrinsic distance and small-sphere maps]
\label{def:chow-iii-i28-sphere-maps}
\source{chow-et-al-ricci-flow-part-iii:page-0489,chow-et-al-ricci-flow-part-iii:page-0490}
\dcref{appi:eq-1.13,appi:eq-1.14}\group{Distance spheres}
For $0<s\le t<\inj(p)$ let
\[
\varphi_{s,t}=\exp_p\circ(t/s)\circ(\exp_p|_{\overline B(0,s)})^{-1}:S(p,s)\to S(p,t),
\]
and let $\dhat_{S(p,s)}$ be the intrinsic length distance induced by the
subspace metric on $S(p,s)$.  In Euclidean space the corresponding maps are
isometries after scaling by $1/s$ and $1/t$.
\end{definition}

\begin{proposition}[Monotonicity inside the injectivity radius]
\label{prop:chow-iii-i29-small-sphere-monotonicity}
\source{chow-et-al-ricci-flow-part-iii:page-0490,chow-et-al-ricci-flow-part-iii:page-0491,chow-et-al-ricci-flow-part-iii:page-0492}
\uses{def:chow-iii-i28-sphere-maps}
\dcref{appi:1.29}\group{Distance spheres}

If $\sect(g)\ge0$ and $s\le t\le u<\inj(p)$, then
\[
\dhat_{S(p,t)}(\varphi_{s,t}x,\varphi_{s,t}y)
\le\frac ts\dhat_{S(p,s)}(x,y),\qquad
\varphi_{t,u}\circ\varphi_{s,t}=\varphi_{s,u}.
\]
\end{proposition}
\begin{proof}
\sketch
The composition identity is immediate from the exponential definition.  A
Jacobi field $J$ with $J(0)=0$ satisfies Rauch comparison
$|J(t)|/|J(s)|\le t/s$, and integrating this estimate along an intrinsic
minimizing path on $S(p,s)$ gives the inequality.
\end{proof}

\begin{definition}[Set gradient and regular points]
\label{def:chow-iii-i30-i31-gradient-regularity}
\source{chow-et-al-ricci-flow-part-iii:page-0493,chow-et-al-ricci-flow-part-iii:page-0494}
\dcref{appi:1.30,appi:1.31}\group{Critical point theory}
Fix $r(x)=d(x,p)$.  Define
\[
(\nabla_\star r)(x)=\{V\in T_x\M:|V|=1,\ u+r(\exp_x(-uV))=r(x)\ \forall u\in[0,r(x)]\}.
\]
A point is regular if some unit vector makes angle $<\pi/2$ with every element
of $(\nabla_\star r)(x)$; otherwise it is critical.
\end{definition}

\begin{lemma}[Local compactness and openness]
\label{lem:chow-iii-i30-i33-regular-open}
\source{chow-et-al-ricci-flow-part-iii:page-0493,chow-et-al-ricci-flow-part-iii:page-0495}
\uses{def:chow-iii-i30-i31-gradient-regularity}
\dcref{appi:1.30,appi:1.33}\group{Critical point theory}

The set gradient $\nabla_\star r$ is locally compact.  For
$0<\theta\le\pi/2$, the set of $\theta$-regular points and the corresponding
cone of admissible directions are open.
\end{lemma}
\begin{proof}
\sketch
Convergent unit vectors determine convergent minimizing geodesics by continuous
dependence of the exponential map, so the limit remains in $\nabla_\star r$.
The openness statement follows by contradiction and compactness of the set
gradient.
\end{proof}

\begin{lemma}[Characterization of regular points]
\label{lem:chow-iii-i34-i35-regular-characterization}
\source{chow-et-al-ricci-flow-part-iii:page-0495,chow-et-al-ricci-flow-part-iii:page-0496,chow-et-al-ricci-flow-part-iii:page-0497}
\uses{def:chow-iii-i30-i31-gradient-regularity, lem:chow-iii-i30-i33-regular-open}
\dcref{appi:1.34,appi:1.35}\group{Critical point theory}

For $x\ne p$, $x$ is regular iff there exist $V\in S_x^{n-1}$, $c>0$, and
$\varepsilon>0$ such that
\[
r(\exp_x(sV))-r(x)\ge cs\qquad(0\le s<\varepsilon).
\]
Equivalently, every sufficiently short unit-speed path initially tangent to a
suitable $V$ increases $r$ at a uniform positive rate.
\end{lemma}

\begin{lemma}[Maximum angle tends to zero]
\label{lem:chow-iii-i37-angle-decay}
\source{chow-et-al-ricci-flow-part-iii:page-0497,chow-et-al-ricci-flow-part-iii:page-0498,chow-et-al-ricci-flow-part-iii:page-0499}
\uses{def:chow-iii-i30-i31-gradient-regularity, lem:chow-iii-i34-i35-regular-characterization}
\dcref{appi:1.37}\group{Critical point theory}

If $\sect(g)\ge0$, the maximum angle
$\operatorname{ang}_p(x)=\max\{\angle(V,W):V,W\in(\nabla_\star r)(x)\}$
satisfies $\operatorname{ang}_p(x)\to0$ as $d(x,p)\to\infty$.  Thus all
sufficiently distant points are $\theta$-regular for every fixed
$0<\theta\le\pi/2$.
\end{lemma}

\begin{definition}[Lipschitz hypersurface]
\label{def:chow-iii-i39-lipschitz-hypersurface}
\source{chow-et-al-ricci-flow-part-iii:page-0500}
\dcref{appi:1.39}\group{Distance spheres}
A subset of a smooth manifold is a Lipschitz hypersurface when in every local
coordinate chart it is the graph of a Lipschitz function on an open subset of
$\R^{n-1}$.
\end{definition}

\begin{lemma}[Completeness of closed Lipschitz hypersurfaces]
\label{lem:chow-iii-i40-lipschitz-complete}
\source{chow-et-al-ricci-flow-part-iii:page-0500}
\uses{def:chow-iii-i39-lipschitz-hypersurface}
\dcref{appi:1.40}\group{Distance spheres}

If $\N$ is connected, closed, and Lipschitz in a complete Riemannian
manifold, then its intrinsic length metric is complete.
\end{lemma}
\begin{proof}
\sketch
Constant-speed minimizing sequences of paths are uniformly Lipschitz.  Arzela--
Ascoli gives a limiting path in the closed set, and lower semicontinuity of
length shows that it realizes the intrinsic distance.
\end{proof}

\begin{proposition}[Large distance spheres]
\label{prop:chow-iii-i41-large-distance-spheres}
\source{chow-et-al-ricci-flow-part-iii:page-0501,chow-et-al-ricci-flow-part-iii:page-0502,chow-et-al-ricci-flow-part-iii:page-0503}
\uses{lem:chow-iii-i37-angle-decay, lem:chow-iii-i40-lipschitz-complete}
\dcref{appi:1.41}\group{Distance spheres}

For all sufficiently large $s$, $S(p,s)$ is a Lipschitz hypersurface, its
intrinsic length space is complete, and the number of its components equals the
number of topological ends of $\M$.
\end{proposition}

\section{Approximate Busemann--Feller theorem}

\begin{lemma}[Tubular neighborhood projection]
\label{lem:chow-iii-i42-tubular-projection}
\source{chow-et-al-ricci-flow-part-iii:page-0504,chow-et-al-ricci-flow-part-iii:page-0505,chow-et-al-ricci-flow-part-iii:page-0506}
\dcref{appi:1.42}\group{Tubular neighborhoods}
Let $\mathcal H$ be a complete two-sided hypersurface with
$\sec(g)\le\kappa^2$ and $\mathrm{II}_{\mathcal H}\ge-\Lambda I$.  There is
\[
\ell(\Lambda,\kappa)=\kappa^{-1}\bigl(\pi-\cot^{-1}(-\kappa^{-1}\Lambda)\bigr)>0
\]
such that the normal exponential map is a diffeomorphism on the positive
tube of radius $\ell(\Lambda,\kappa)$; nearest-point projection
$\pi:\mathcal N^+_{\ell}(\mathcal H)\to\mathcal H$ is smooth and unique.
\end{lemma}

\begin{theorem}[Approximate Busemann--Feller projection]
\label{thm:chow-iii-i44-approximate-busemann-feller}
\source{chow-et-al-ricci-flow-part-iii:page-0506,chow-et-al-ricci-flow-part-iii:page-0507,chow-et-al-ricci-flow-part-iii:page-0508}
\uses{lem:chow-iii-i42-tubular-projection}
\dcref{appi:1.43,appi:1.44}\group{Tubular neighborhoods}

For every smooth path $\gamma$ in the positive tube of radius
$\varepsilon\le\ell(\Lambda,\kappa)$,
\[
\mathcal L(\pi\circ\gamma)\le
\left(\cos(\varepsilon\kappa)-\frac\Lambda\kappa\sin(\varepsilon\kappa)\right)^{-1}
\mathcal L(\gamma).
\]
\end{theorem}

\section{Mollified distance and points at infinity}

\begin{lemma}[Mollified distance]
\label{lem:chow-iii-i45-mollified-distance}
\source{chow-et-al-ricci-flow-part-iii:page-0509}
\dcref{appi:1.45}\group{Mollified distance}
For the standard mollification $r_\varepsilon$ of $r=d(\cdot,p)$, one has
$r_\varepsilon\to r$ uniformly on compact sets and
$\nabla r_\varepsilon\to\nabla r$ uniformly on compact subsets of the complement
of the cut locus.
\end{lemma}

\begin{lemma}[Derivative estimates for mollified distance]
\label{lem:chow-iii-i47-mollified-derivatives}
\source{chow-et-al-ricci-flow-part-iii:page-0510,chow-et-al-ricci-flow-part-iii:page-0511}
\uses{lem:chow-iii-i45-mollified-distance, lem:chow-iii-i37-angle-decay}
\dcref{appi:1.47}\group{Mollified distance}

On a complete noncompact manifold with $\sect(g)\ge0$, for each compact ball
there are errors $\kappa_i(\varepsilon)\to0$ such that
\[
1-\kappa_1-\theta_1(r)\le|\nabla r_\varepsilon|\le1+\kappa_1,
\quad |\nabla r_\varepsilon|\ge1-\kappa_2
\]
away from the cut locus, and
$\nabla^2r_\varepsilon\le(1+\kappa_3)/r_\varepsilon$.
\end{lemma}

\begin{example}[Mollification and intrinsic-distance distortion]
\label{ex:chow-iii-i46-i49-mollification}
\source{chow-et-al-ricci-flow-part-iii:page-0510,chow-et-al-ricci-flow-part-iii:page-0511,chow-et-al-ricci-flow-part-iii:page-0512}
\dcref{appi:1.46,appi:1.49}\group{Mollified distance}
On a sphere, a mollified distance has zero gradient at the antipode.  Smooth or
Lipschitz hypersurfaces may converge uniformly to a plane while their intrinsic
lengths do not converge, so intrinsic-distance control is an additional issue.
\end{example}

\begin{claim}[Presumed distance-nonincreasing maps between large spheres]
\label{clm:chow-iii-i50-large-sphere-maps}
\source{chow-et-al-ricci-flow-part-iii:page-0512,chow-et-al-ricci-flow-part-iii:page-0513}
\uses{prop:chow-iii-i29-small-sphere-monotonicity}
\dcref{appi:1.50}\group{Points at infinity}

For sufficiently large $s_0(p)$ and $t\ge s\ge s_0(p)$, there should exist maps
\[
\phi_{s,t}:S(p,s)\to S(p,t)
\]
that are nonincreasing for the rescaled intrinsic metrics, satisfy
$\phi_{t,u}\circ\phi_{s,t}=\phi_{s,u}$, and send $\alpha(s)$ to $\alpha(t)$
for each ray $\alpha$ from $p$.  This is the source's Presumed Theorem I.50.
\end{claim}

\begin{definition}[Equivalence of rays]
\label{def:chow-iii-i51-ray-equivalence}
\source{chow-et-al-ricci-flow-part-iii:page-0513}
\dcref{appi:1.51}\group{Points at infinity}
Two rays are equivalent, $\alpha\sim\beta$, when
$d(\alpha(s),\beta(s))/s\to0$.  Let $[\alpha]$ denote the equivalence class.
\end{definition}

\begin{example}[Euclidean ideal boundary]
\label{ex:chow-iii-i51-euclidean-boundary}
\source{chow-et-al-ricci-flow-part-iii:page-0513}
\dcref{appi:ex-1.51}\group{Points at infinity}
For $\M=\mathbb R^n$, the quotient of rays by $\sim$ is $\mathbb S^{n-1}$: rays are
equivalent exactly when they have the same direction.
\end{example}

\begin{lemma}[Rays based at any point]
\label{lem:chow-iii-i52-rays-at-basepoint}
\source{chow-et-al-ricci-flow-part-iii:page-0513}
\uses{def:chow-iii-i51-ray-equivalence}
\dcref{appi:1.52}\group{Points at infinity}

For every ray $\alpha$ and every $p\in\M$, there is a ray $\alpha_p\in\Ray_{\M}(p)$
with $\alpha_p\sim\alpha$.
\end{lemma}

\begin{definition}[Ray-space distance]
\label{def:chow-iii-i53-ray-distance}
\source{chow-et-al-ricci-flow-part-iii:page-0513}
\uses{def:chow-iii-i51-ray-equivalence}
\dcref{appi:eq-1.62}\group{Points at infinity}

Whenever the limit exists, define
\[
d_R([\alpha],[\beta])=lim_{s\to\infty}
\frac{\dhat_{S(p,s)}(\alpha\cap S(p,s),\beta\cap S(p,s))}{s}.
\]
\end{definition}

\begin{lemma}[Well-definedness of ray-space distance]
\label{lem:chow-iii-i53-ray-distance-well-defined}
\source{chow-et-al-ricci-flow-part-iii:page-0514,chow-et-al-ricci-flow-part-iii:page-0515}
\uses{def:chow-iii-i53-ray-distance, lem:chow-iii-i52-rays-at-basepoint, clm:chow-iii-i50-large-sphere-maps}
\dcref{appi:1.53}\group{Points at infinity}

For nonnegative sectional curvature, $d_R$ is independent of representatives;
$d_R=0$ exactly for equivalent rays.  Assuming Presumed Theorem I.50, the
limit defining $d_R$ exists.
\end{lemma}

\begin{claim}[Presumed compact space of points at infinity]
\label{clm:chow-iii-i54-points-at-infinity}
\source{chow-et-al-ricci-flow-part-iii:page-0515,chow-et-al-ricci-flow-part-iii:page-0516}
\uses{lem:chow-iii-i53-ray-distance-well-defined}
\dcref{appi:1.54}\group{Points at infinity}

Assuming Presumed Theorem I.50, $(\Ray_{\M}/\sim,d_R)$ is a compact metric space of
finite diameter, called the space of points at infinity.
\end{claim}

\begin{remark}[Tits metric]
\label{rem:chow-iii-i55-tits-metric}
\source{chow-et-al-ricci-flow-part-iii:page-0516}
\dcref{appi:1.55}\group{Points at infinity}
The analogous quotient metric for simply connected complete manifolds of
nonpositive curvature is the Tits metric; its metric cone is the Tits cone.
\end{remark}
